\subsection{Readout error mitigation}
\label{sec:readout}

Readout errors map the ideal distribution to the measured one,
\begin{equation}
\vec{P}_{\rm meas}=M\vec{P}_{\rm ideal},
\end{equation}
where $M$ is the assignment matrix with $M_{ab}=P(a|b)$.
Its columns are calibrated by preparation: preparing $|0\rangle$
gives the first column, preparing $|1\rangle$ the second,
\begin{equation}
M=\begin{bmatrix} P(0|0) & P(0|1) \\ P(1|0) & P(1|1) \end{bmatrix},
\end{equation}
where the diagonals read correct and the off-diagonals misread.
Mitigation inverts the map,
\begin{equation}
\vec{P}_{\rm mit}=M^{-1}\vec{P}_{\rm meas},
\end{equation}
where $\vec{P}_{\rm mit}$ is then projected onto the simplex by
truncating negatives to zero and renormalizing. In other words, the
raw counts are unfolded with the calibrated confusion matrix before any
observable is evaluated, following the tensor-product readout mitigation of
Bravyi~\textit{et al.}~\cite{bravyi2021mitigating}.
Assuming uncorrelated per-qubit errors, each qubit is corrected
independently at a calibration cost of $2$ circuits per qubit, linear in
system size; on hardware we recalibrate per batch against drift, with the
platform's built-in readout correction switched off.
The uncorrelated assumption keeps the overhead linear but cannot capture
crosstalk; correlated readout errors, if present, survive mitigation.
This tensor-product inversion sits within a broader mitigation
landscape. Probabilistic error cancellation expands noisy gates in a
quasiprobability basis and resamples
accordingly~\cite{temme2017error}, now extended to dynamic circuits with
mid-circuit measurement~\cite{gupta2024probabilistic} and unified with
noise scaling~\cite{mari2021extending}; zero-noise extrapolation instead
fits $E(\lambda)$ at boosted noise strengths and continues to
$E(0)$---first demonstrated on superconducting
hardware in Ref.~\cite{kandala2019error}---with optimized Richardson
protocols~\cite{krebsbach2022optimization}, inverted-circuit noise
estimation~\cite{koenig2024inverted}, and active error suppression by
artificial noise boosting inside variational
loops~\cite{li2017efficient}---though time-correlated noise is
systematically harder to extrapolate
away~\cite{schultz2022analyzing}. Matrix-free iterative unfolding scales
the same assignment-matrix idea to dozens of qubits without ever forming
$M$~\cite{nation2021scalable}; our per-qubit inversion is its minimal,
drift-robust endpoint.
